判别模型直接学习输入到输出的关系，生成式模型则先描述类别、潜在状态与观测如何共同产生。后一种路线需要更多假设，却能自然处理缺失标签、软归属、序列状态和不确定性。

本单元不是模型的陈列，而是一条逐步隐藏信息的路径：

\begin{enumerate}
\def\labelenumi{\arabic{enumi}.}
\tightlist
\item
  生成式分类从 Naive Bayes 到 GDA：标签可见时，先学习各类怎样生成特征，再用 Bayes 公式反推类别。
\item
  Gaussian 混合模型让类别成为隐变量：拿走标签后，每个样本属于哪个成分也需要推断，硬聚类变成 posterior responsibility。
\item
  EM 用下界交替完成隐变量学习：隐变量未知使似然中出现难以拆开的对数求和，EM 用当前后验和参数更新交替推进。
\item
  HMM 把隐变量排列成时间链：让隐状态不再彼此独立，而沿时间转移；动态规划由此成为推断核心。
\end{enumerate}

学习时要始终区分三层：模型假设规定了什么分布，推断算法近似或计算了什么量，数据是否足以识别我们赋予隐变量的现实含义。一个算法收敛，并不能替代后两项判断。
